\begin{section}{The Real Field}

There exists an ordered field $\R$ which has the \hyperref[def:LeastUpperBoundProperty]{least upper bound property}.
Moreover, $\R$ contains $\Q$ as a subfield, i.e., $\Q \subseteq \R$ and the operations of addition and multiplication
on $\R$ coincide with those on $\Q$.
\\[3mm]
We can construct $\R$ from $\Q$. But for now, we will assume the existence of $\R$ and its properties.

\bigskip

\begin{theorem} \label{thm:ArchimedeanProperty}
    If $x, y \in \R$ and $x > 0$, then there exists some $n \in \N$ such that $nx > y$.
    This theorem is called the \textbf{Archimedean property} of $\R$.
\end{theorem}

\begin{proof}
    Suppose $x, y \in \R$ and $x > 0$. Let $E \coloneqq \setbuilder{nx}{nx \leq y, \; n \in \N}$.
    If $E = \emptyset$, which implies that $x > y$, $n = 1$ satisfies $nx > y$.
    Otherwise, $E$ is nonempty and bounded above by $y$.
    \\[3mm]
    By the least upper bound property, there exists $\alpha = \sup E \in \R$.
    Then $\alpha - x$ is not an upper bound of $E$. Therefore, there exists some
    $m \in \N$ such that $mx > \alpha - x \implies (m + 1)x > \alpha$. Then $(m + 1)x \notin E$,
    which implies that $(m + 1)x > y$. Therefore, we can take $n = m + 1$.
\end{proof}

\bigskip

\begin{corollary} \label{cor:ArchimedeanPropertyCorollary}
    For any $\epsilon > 0$, there exists some $n \in \N$ such that $b^{-n} < \epsilon$ for any $b \in \R$ with $b > 1$.
\end{corollary}

\begin{proof}
    For all $n \in \N$ and $x \in \R$ with $x \geq -1$, we have $(1 + x)^{n} \geq 1 + nx$. Let $a \in \R^{+}$ be given.
    By \Cref{thm:ArchimedeanProperty}, there exists some $n \in \N$ such that $na > \dfrac{1}{\epsilon}$. Then we have
    \[
    (1 + a)^{n} \geq 1 + na > na > \frac{1}{\epsilon} \implies (1 + a)^{-n} < \epsilon.
    \]
    Since $b > 1$, we can take $a = b - 1 > 0$ to complete the proof.
\end{proof}

\bigskip

\begin{remark}
    \Cref{thm:BoundedSubsetOfZHasExtremum} supposes that $E$ is bounded in $\Z$. By \Cref{thm:ArchimedeanProperty},
    we can always find an integer $k$ such that $k < x$ or $x < k$ for any $x \in \R$. 
    Therefore, \Cref{thm:BoundedSubsetOfZHasExtremum} also holds when $E$ is bounded in $\R$, which is followed by
    the fact that $E$ is also bounded in $\Z$.
\end{remark}

\bigskip

\begin{definition} \label{def:FloorCeilingFunctions}
    For any $x \in \R$, the \textbf{floor function} $\floor{x}$
    and \textbf{ceiling function} $\ceil{x}$ are defined as follows:
    \[
    \floor{x} \coloneqq \max\setbuilder{k \in \Z}{k \leq x}, \quad
    \ceil{x} \coloneqq \min\setbuilder{k \in \Z}{k \geq x}
    \]
\end{definition}

\bigskip

\begin{theorem} \label{thm:FloorCeilingFunctionsExistenceAndUniqueness}
    Let $x \in \R$ be given. Then $\floor{x}$ and $\ceil{x}$ exist and are unique.
\end{theorem}

\begin{proof}
    Let $E \coloneqq \setbuilder{k \in \Z}{k \leq x}$. If $x \geq 0$, then $-1 \in E$. Otherwise, by
    \Cref{thm:ArchimedeanProperty}, there exists some $n \in \N$ such that $n > -x$, which implies that $-n \in E$.
    Then, by \Cref{thm:BoundedSubsetOfZHasExtremum}, $E$ has a unique greatest element.
    \\[6mm]
    Let $F \coloneqq \setbuilder{k \in \Z}{k \geq x}$. By \Cref{thm:ArchimedeanProperty}, there exists some $n \in \N$
    such that $n > x$, which implies that $n \in F$. Then, by \Cref{thm:BoundedSubsetOfZHasExtremum},
    $F$ has a unique least element.
\end{proof}

\bigskip

\begin{theorem} \label{thm:QIsDenseInR}
    If $x, y \in \R$ and $x < y$, then there exists some $p \in \Q$ such that $x < p < y$.
    This theorem states that $\Q$ is \textbf{dense} in $\R$.
\end{theorem}

\begin{proof}
    Let $\epsilon = y - x > 0$. By \Cref{thm:ArchimedeanProperty}, there exists some $n \in \N$
    such that $\dfrac{1}{n} < \epsilon$. Let $S \coloneqq \setbuilder*{k \in \Z}{\dfrac{k}{n} < x}$.
    Since $n\floor{x} - 1 \in S$, $S$ is nonempty and bounded above by $nx$.
    By \Cref{thm:BoundedSubsetOfZHasExtremum}, $S$ has a greatest element $m = \max{S}$.
    \[
    x < \frac{m + 1}{n} < x + \frac{1}{n} < x + (y - x) = y.
    \]
    Therefore, we can take $p = \dfrac{m + 1}{n} \in \Q$.
\end{proof}

\bigskip

\begin{theorem} \label{thm:ExistenceOfNthRoots}
    For each $x \in \R^{+}$ and $n \in \N$, there exists a unique $y \in \R^{+}$ such that
    $y^{n} = x$. We denote $y$ by $\sqrt[n]{x}$ or $x^{\frac{1}{n}}$.
\end{theorem}

\begin{proof}
    Let $E \coloneqq \setbuilder{t \in \R^{+}}{t^{n} < x}$. Then $E \neq \emptyset$ since
    $2^{-1} \in E$ if $x\geq 1$, and $x^{2} \in E$ otherwise. Moreover, $E$ is bounded above by $\max\set{x, 1}$.
    By the least upper bound property, there exists $\alpha = \sup E \in \R^{+}$. We will show that $\alpha^{n} = x$.
    \\[6mm]
    For any $a, b \in \R^{+}$ with $a < b$, we have
    \[
    b^{n} - a^{n} = (b - a)\sum_{k = 0}^{n - 1} b^{n - 1 - k} a^{k} < (b - a)nb^{n - 1}.
    \]
    Assume that $\alpha^{n} < x$. By \Cref{thm:QIsDenseInR}, we can choose $h \in \R^{+}$ such that
    \[
    h < \min\set*{1, \; \frac{x - \alpha^{n}}{n(\alpha + 1)^{n - 1}}}.
    \]
    Put $a = \alpha$ and $b = \alpha + h$. Then we have
    \begin{align*}
        &(\alpha + h)^{n} - \alpha^{n} < hn(\alpha + h)^{n - 1} < hn(\alpha + 1)^{n - 1} \\
        \implies &(\alpha + h)^{n} < \alpha^{n} + hn(\alpha + 1)^{n - 1} < \alpha^{n} + (x - \alpha^{n}) = x.
    \end{align*}
    This shows that $\alpha + h \in E$, which contradicts the fact that $\alpha = \sup E$.
    \\[6mm]
    Assume that $\alpha^{n} > x$. By \Cref{thm:QIsDenseInR}, we can choose $h \in \R^{+}$ such that
    \[
    h < \min\set*{1, \; \alpha, \; \frac{\alpha^{n} - x}{n(\alpha + 1)^{n - 1}}}.
    \]
    Put $a = \alpha - h$ and $b = \alpha$. Then we have
    \begin{align*}
        &\alpha^{n} - (\alpha - h)^{n} < hn\alpha^{n - 1} < hn(\alpha + 1)^{n - 1} \\
        \implies &(\alpha - h)^{n} > \alpha^{n} - hn(\alpha + 1)^{n - 1} > \alpha^{n} - (\alpha^{n} - x) = x.
    \end{align*}
    This shows that $\alpha - h$ is an upper bound of $E$, which contradicts the fact that $\alpha = \sup E$.
    \\[6mm]
    Therefore, we have $\alpha^{n} = x$, which guarantees the \emph{existence} of $y$. For any $\alpha, \beta \in \R^{+}$,
    $\alpha \neq \beta \implies \alpha^{n} \neq \beta^{n}$. Thus the \emph{uniqueness} of $y$ is guaranteed.
\end{proof}

\bigskip

\begin{definition} \label{def:GreedyBaseExpansionOfRealNumbers}
    Let $y \in [0, 1)$ and $b \in \N$ be given with $b \geq 2$. A sequence $\set{a_{n}}_{n = 1}^{\infty}$
    is called the \textbf{greedy base $b$ expansion} of $y$ if for each $n \in \N$, we have
    \[
    0 \leq y - \sum_{k = 1}^{n} \frac{a_{k}}{b^{k}} < \frac{1}{b^{n}}
    \]
    where $a_{k} \in S_{b} \coloneqq \set{0, 1, 2, \dots, b - 1}$ for all $k \in \N$.
    Especially, if $b = 10$, then the sequence is called the \textbf{greedy decimal expansion} of $y$.
\end{definition}

\bigskip

\begin{theorem} \label{thm:ExistenceAndUniquenessOfGreedyBaseExpansion}
    For any $y \in [0, 1)$, the greedy base $b$ expansion of $y$ exists and is unique. 
\end{theorem}

\begin{proof}
    We define sequences $\set{t_{n}}_{n = 1}^{\infty}$, $\set{T_{n}}_{n = 1}^{\infty}$
    and $\set{a_{n}}_{n = 1}^{\infty}$ as follows:
    \begin{itemize}
        \item $t_{1} = by$, $T_{1} \coloneqq \setbuilder{t \in S_{b}}{t \leq t_{1}}$, $a_{1} = \max T_{1}$.
        \item $t_{2} = b^{2}y - ba_{1}$, $T_{2} \coloneqq \setbuilder*{t \in S_{b}}{t \leq t_{2}}$, $a_{2} = \max T_{2}$.
        \item[] $\vdots$ 
        \item $t_{n} = b^{n}y - \displaystyle \sum_{k = 1}^{n - 1} a_{k}b^{n - k}$,
        $T_{n} \coloneqq \setbuilder*{t \in S_{b}}{t \leq t_{n}}$, $a_{n} = \max T_{n}$
        
        for all $n \in \N$ with $n \geq 2$,
    \end{itemize}
    Suppose that $t_{n} \geq 0$ for some $n \in \N$. Since $0 \in T_{n}$ and $T_{n}$ is bounded above, $T_{n}$ has
    a greatest element by \Cref{thm:BoundedSubsetOfZHasExtremum}, which implies that $\set{a_{n}}_{n = 1}^{\infty}$
    is well-defined. Let $P(n):$ $t_{n} \geq 0$ and $a_{n}$ is well-defined.
    \begin{itemize}
        \item \textbf{Base case:} $P(1)$ is true since $t_{1} = by \geq 0$.
        \item \textbf{Inductive step:} Suppose that $P(n)$ is true for some $n \in \N$. Then we have
        \begin{align*}
            t_{n + 1} &= b^{n + 1}y - \sum_{k = 1}^{n} a_{k}b^{n + 1 - k} \\
            &= b\paren*{b^{n}y - \sum_{k = 1}^{n} a_{k}b^{n - k}} \\
            &= b\paren*{b^{n}y - \sum_{k = 1}^{n - 1} a_{k}b^{n - k} - a_{n}} \\
            &= b\paren{t_{n} - a_{n}} \geq 0.
        \end{align*}
        Therefore, $P(n + 1)$ is true whenever $P(n)$ is true.
    \end{itemize}
    By the principle of mathematical induction, $P(n)$ is true for all $n \in \N$.
    Then, by the definition of $a_{n}$, we have
    \begin{align*}
        &a_{n} \leq b^{n}y - \sum_{k = 1}^{n - 1} a_{k}b^{n - k} < a_{n} + 1 \\
        \implies &0 \leq b^{n}y - \sum_{k = 1}^{n} a_{k}b^{n - k} < 1 \\
        \implies &0 \leq y - \sum_{k = 1}^{n} \frac{a_{k}}{b^{k}} < \frac{1}{b^{n}}
    \end{align*}
    for each $n \in \N$. Therefore, the \emph{existence} of the \emph{greedy base $b$ expansion} of $y$ is guaranteed.
    \\[6mm]
    Next, suppose that there exist two sequences $\set{a_{n}}_{n = 1}^{\infty}$
    and $\set{c_{n}}_{n = 1}^{\infty}$ such that for each $n \in \N$, we have